\chapter{The Renormalization Battle}

The Wilson line of \cref{eq:quasipdf} is a self-energy divergent object:
beyond the usual logarithms it carries a \emph{linear} UV divergence,
proportional to its length $z$. Whether and how this can be tamed decided
the fate of lattice parton physics. This chapter follows the battle in
seven papers: from the veteran renormalization scheme that made lattice
matrix elements respectable, through the proofs of renormalizability, to
the practical schemes (ratio/RI/MOM hybrids, auxiliary fields,
self-renormalization) in daily use inside this library.

% ----------------------------------------------------------------------
\section{RI/MOM: The Veteran Scheme}
\papercard{G.\ Martinelli, C.\ Pittori, C.\ T.\ Sachrajda, M.\ Testa,
A.\ Vladikas}
{A general method for nonperturbative renormalization of lattice operators}
{Nucl.\ Phys.\ B \textbf{445}, 81--108 (1995)}
{arXiv:hep-lat/9411010 | DOI: 10.1016/0550-3213(95)00126-D}
\whyselected{The founding scheme of nonperturbative operator
renormalization ($\sim$8 explicit citations plus universal implicit use):
every RI/MOM or SMOM renormalized matrix element in the library descends
from it.}
\physics{Renormalize an operator $O$ by imposing a renormalization
condition on its matrix element in a family of off-shell Landau-gauge
momentum-subtraction states,
\begin{equation}
  Z_O(\mu,a)\,\Lambda_O(p)\big|_{p^2=\mu^2} = \Lambda_O^{tree}(p),
  \qquad p^2 \gg \Lambda_{\rm QCD}^2 ,
\end{equation}
computed in lattice perturbation-free simulations. The scheme is
renormalization-group closed with computable conversion factors to
$\overline{\rm MS}$, and --- crucially for later chapters --- extends
verbatim to bilocal operators, where its momentum-space amputated vertex
became the diagnostic instrument for the linear divergence of Wilson
lines.}
\impact{All early quasi-PDF extractions in the library (Liu et al.,
helicity at physical mass, isovector systematic study) renormalize
RI/MOM-style; the TMD papers document where it fails (Chapter~6 context),
which is itself part of this paper's legacy: you can only diagnose a
scheme's limits because the scheme exists.}
\cardend

% ----------------------------------------------------------------------
\section{Renormalizability Proved}
\papercard{T.\ Ishikawa, Y.-Q.\ Ma, J.-W.\ Qiu, S.\ Yoshida}
{Renormalizability of quasiparton distribution functions}
{Phys.\ Rev.\ D \textbf{96}, 094019 (2017)}
{arXiv:1707.03107 | DOI: 10.1103/PhysRevD.96.094019}
\whyselected{First general proof that quasi-PDF operators are
renormalizable ($\sim$13/99): the license without which the program might
have been a numerical artifact.}
\physics{Analyzing all one-loop diagrams of the bilocal operator, the
paper shows that every power divergence localizes in the cusp-less part
of the Wilson-line self-energy --- i.e.\ the operator renormalizes
multiplicatively up to mixing among operators of equal dimension with
same quantum numbers, and no new nonlocal counterterm is required.
Combined with a gauge-invariant UV regulator argument, the result
guarantees that a continuum limit of suitably renormalized quasi-PDFs
exists. The companion PRL by Ma--Qiu (arXiv:1404.6860 lineage; ``good
lattice cross sections'') framed the same logic as a general strategy:
compute at short distances where the lattice is reliable, match to the
long-distance observable.}
\impact{Every renormalization proposal in this library --- ratio schemes,
auxiliary-field subtraction, hybrid schemes --- cites this proof as its
theoretical warrant.}
\cardend

% ----------------------------------------------------------------------
\section{Taming the Linear Divergence}
\papercard{J.-W.\ Chen, X.\ Ji, J.-H.\ Zhang}
{Improved quasi parton distribution through Wilson line renormalization}
{Nucl.\ Phys.\ B \textbf{915}, 1--9 (2017)}
{arXiv:1609.08102 | DOI: 10.1016/j.nuclphysb.2016.12.011}
\whyselected{Origin of the auxiliary-field (mass counterterm) cure for the
linear divergence ($\sim$11/99): cited by essentially every
renormalization-oriented article in the corpus.}
\physics{Introduce an auxiliary field $\phi$ whose quadratic term carries
a mass $m_0$; the open Wilson line re-expresses as expectation values in
an enlarged theory where the linear divergence appears solely as the
self-energy of $\phi$, absorbed by
\begin{equation}
  m_0^2(z,a)=m^2 + \frac{c_1}{a} + \cdots .
\end{equation}
After subtraction the improved quasi-distribution contains only logarithmic
divergences at any order in perturbation theory, and admits standard
nonperturbative renormalization. The paper also computed the first
lattice-regularized one-loop matching kernel, exposing how discretization
artifacts enter $K$.}
\impact{This construction is the seed of Green--Jansen--Steffens'
implementation and of the hybrid schemes below; conceptually it converts
``subtracting a linear divergence'' from a numerics problem into a mass
renormalization problem --- solvable by fitting $m_0(z)$ across $z$.}
\cardend

% ----------------------------------------------------------------------
\section{The Renormalization Program Formalized}
\papercard{X.\ Ji, J.-H.\ Zhang, Y.\ Zhao}
{Renormalization in large momentum effective theory of parton physics}
{Phys.\ Rev.\ Lett.\ \textbf{120}, 112001 (2018)}
{arXiv:1706.08962 | DOI: 10.1103/PhysRevLett.120.112001}
\whyselected{The programmatic statement of LaMET renormalization
($\sim$13/99): the reference for auxiliary-field formalism and for the
Wilson-line anomalous dimension entering matching at NNLO.}
\physics{The letter systematizes the two-step logic now used everywhere:
(i)~renormalize the quasi-operator nonperturbatively within a chosen
scheme (RI/MOM, ratio, or auxiliary-field mass), rendering the matrix
element finite in the continuum limit; (ii)~convert to $\overline{\rm MS}$
and match to the light-cone distribution with kernels that include the
cusp logarithm $\ln(P^z/\mu)$. It also computes the one-loop conversion
factors and demonstrates scheme independence of the final light-cone PDF
order by order.}
\impact{The library's hybrid-renormalization applications (transversity,
Boer--Mulders, gluon PDFs) cite this letter as their definitional
framework; its notation is the de facto standard across LP3 papers.}
\cardend

% ----------------------------------------------------------------------
\section{Auxiliary Fields Implemented}
\papercard{J.\ Green, K.\ Jansen, F.\ Steffens}
{Nonperturbative renormalization of nonlocal quark bilinears from lattice QCD}
{Phys.\ Rev.\ Lett.\ \textbf{121}, 022004 (2018)}
{arXiv:1707.07152 | DOI: 10.1103/PhysRevLett.121.022004}
\whyselected{First numerical demonstration that the auxiliary-field scheme
works on realistic ensembles ($\sim$14/99): the bridge between the
Chen--Ji--Zhang idea and everyday practice.}
\physics{On $N_f=2+1$ domain-wall ensembles the paper measures the
quasi-PDF operator's Green function, fits the linearly divergent
$\phi$-field mass $m_0(z,a)$ over a range of separations and spacings,
subtracts it, and shows the residual matrix element scales
logarithmically toward a finite continuum limit --- then supplies the
conversion factors to $\overline{\rm MS}$. The fit-based extraction of a
power-divergent quantity, with its correlated systematics, became the
template that self-renormalization (below) would generalize.}
\impact{Within the library, the ETMC and LPC lines quote this paper when
justifying hybrid short-distance treatments; its error-analysis patterns
reappear in every modern extraction.}
\cardend

% ----------------------------------------------------------------------
\section{The Complete Prescription}
\papercard{C.\ Alexandrou, K.\ Cichy, M.\ Constantinou, K.\ Hadjiyiannakou,
K.\ Jansen, H.\ Panagopoulos, F.\ Steffens}
{A complete non-perturbative renormalization prescription for quasi-PDFs}
{Nucl.\ Phys.\ B \textbf{923}, 394--415 (2017)}
{arXiv:1706.00265 | DOI: 10.1016/j.nuclphysb.2017.08.010}
\whyselected{The most complete single-scheme treatment ($\sim$10/99): the
ETMC school's standard, shared by many library papers including the
twisted-mass gluon PDF.}
\physics{For unpolarized, helicity and transversity quasi-PDFs the paper
constructs the full RI$'$ prescription: renormalization conditions for
all relevant tensor structures, subtraction of the linear divergence via
the $z$-dependent mass term, inclusion of the mixing with the scalar
operator for twist-3 structures, and perturbative conversion factors
$\mathrm{RI}'\to\overline{\rm MS}$ at one loop. The result is a recipe
that a simulation can follow end-to-end --- matrix element, $Z$ factors,
matching, extrapolations --- which explains its role as the field's
practical benchmark.}
\impact{In this library it anchors the ETMC proton gluon-PDF calculation
and provides the comparison point against hybrid-ratio schemes; the
gauge-fixing-precision paper quantifies exactly how much care its
Landau-gauge step demands.}
\cardend

% ----------------------------------------------------------------------
\section{Gluon Operators Renormalize Multiplicatively}
\papercard{Z.-Y.\ Li, Y.-Q.\ Ma, J.-W.\ Qiu}
{Multiplicative renormalizability of operators defining quasiparton distributions}
{Phys.\ Rev.\ Lett.\ \textbf{122}, 062002 (2019)}
{arXiv:1809.01836 | DOI: 10.1103/PhysRevLett.122.062002}
\whyselected{The rigorous basis for gluon quasi-PDFs ($\sim$8/99): quoted
by nearly every gluon-PDF article of the library.}
\physics{Classifying all gauge-invariant bilocal adjoint operators with a
Wilson line, the paper proves --- using a gauge-invariant UV regulator so
that only genuine UV singularities appear --- that the full set of 36
gluon quasi-parton operators has localized UV divergences and renormalizes
\emph{multiplicatively}: each operator maps onto a finite combination
fixed by symmetry, with no operator mixing beyond the trivial ones. This
completes, for gluons, what Ishikawa et al.\ proved for quarks, and
certifies that lattice gluon quasi-PDFs define continuum quantities.}
\impact{Together with Zhang et al.'s matching analysis (next chapter),
this theorem is why ``first gluon PDF from the lattice'' papers can exist
at all. The library's gluon-quasi-PDF line cites it at every
methodological juncture.}
\cardend
